\chapter{Quantum Circuits and Programming}
\label{ch:circuits}

Quantum algorithms are expressed as \emph{circuits}: wires represent qubits, time flows left to right, boxes are gates, control dots mark controlled operations, and measurement symbols denote readout. This chapter formalizes the circuit model, walks through Bell-state synthesis, surveys gate sets and OpenQASM, and discusses depth and resource costs---connecting mathematics to the software tools used on real hardware.

\section{The Circuit Model}
\label{sec:circuit-model}

\subsection{Unitary evolution}

An $n$-qubit circuit transforms $\ket{\psi_{\mathrm{in}}}$ to
\begin{equation}
\ket{\psi_{\mathrm{out}}}=U_k\cdots U_2 U_1\ket{\psi_{\mathrm{in}}},
\label{eq:circuit-evolution}
\end{equation}
where each $U_i$ acts unitarily on some subset of qubits (often tensored with identity on others). Measurement at the end samples classical outcomes from Born probabilities. Circuits are \emph{acyclic}: no feedback unless classical control is explicitly modeled.

\subsection{Diagrammatic conventions}

Wires index qubits from top to bottom; time flows left to right. A solid control dot on one wire with $\oplus$ on another denotes CNOT. Single-qubit boxes label gates ($H$, $T$, etc.). Double lines occasionally denote classical bits conditioned on measurement outcomes.

\subsection{Depth, width, and gate count}

\begin{itemize}[noitemsep]
\item \textbf{Width:} number of qubits $n$.
\item \textbf{Gate count:} total gates applied.
\item \textbf{Depth:} number of sequential \emph{layers} of gates that can run in parallel on disjoint qubits.
\end{itemize}
Depth determines latency on hardware with sequential gate calibration; gate count relates to total error budget. On superconducting devices, two-qubit gates along the same coupling edge often cannot execute simultaneously.

\begin{example}
Two CNOTs sharing a control qubit cannot run in parallel if they use the same control line in conflicting ways; two single-qubit gates on different qubits \emph{can} occupy one depth layer. A circuit with $n$ sequential CNOTs on a chain has depth $\Omega(n)$ even if gate count is also $n$.
\end{example}

\subsection{Measurement and classical control}

Mid-circuit measurement feeds classical registers that condition later gates (\emph{dynamic circuits}). This extends~\eqref{eq:circuit-evolution} to conditional operations $U_k(c_1,\ldots,c_m)$ where $c_j$ are prior measurement outcomes---essential for teleportation and error correction.

\section{Building Bell States}
\label{sec:bell-circuit}

\subsection{Standard construction}

The canonical Bell circuit prepares $\ket{\Phi^+}$ from $\ket{00}$:
\begin{equation}
\ket{00}\xrightarrow{H_0}\xrightarrow{\mathrm{CNOT}_{0,1}}\ket{\Phi^+}
=\frac{\ket{00}+\ket{11}}{\sqrt{2}}.
\label{eq:bell-circuit}
\end{equation}
Here $H_0$ denotes Hadamard on qubit~0; $\mathrm{CNOT}_{0,1}$ uses qubit~0 as control and qubit~1 as target.

\subsection{Step-by-step state update}

\begin{align*}
\ket{00} &\xrightarrow{H\otimes I} \frac{1}{\sqrt{2}}(\ket{00}+\ket{10}), \\
&\xrightarrow{\mathrm{CNOT}} \frac{1}{\sqrt{2}}(\ket{00}+\ket{11})=\ket{\Phi^+}.
\end{align*}
Other Bell states arise from additional single-qubit Pauli rotations before or after this core.

\begin{labbox}{CB}{Circuit Builder}
Open \texttt{/playground/circuit-builder}. Drag $H$ and CNOT onto two qubits, simulate step-by-step, and read the statevector. Confirm amplitudes match~\eqref{eq:bell-states}.
\end{labbox}

\section{Gate Sets and Native Hardware}
\label{sec:gate-sets}

\subsection{Universal versus native gates}

Theory uses $\{H,T,\mathrm{CNOT}\}$; hardware exposes \emph{native} gates---e.g.\ $R_x,R_y,R_z$ and CNOT on fixed connectivity graphs. \emph{Compilation} transpiles abstract circuits into native pulses with SWAP routing for limited connectivity.

\subsection{Clifford+T and T-count}

Fault-tolerant architectures often separate fast \emph{Clifford} gates ($H,S,\mathrm{CNOT}$) from costly \emph{T} gates, which require magic-state distillation. $T$-count dominates resource estimates for Shor's algorithm.

\section{OpenQASM}
\label{sec:openqasm}

\subsection{Language overview}

Open Quantum Assembly Language (OpenQASM) describes circuits textually. OpenQASM~3.0 adds classical control flow, gate modifiers, and richer typing. A minimal Bell preparation:

\begin{verbatim}
OPENQASM 3.0;
include "stdgates.inc";
qubit[2] q;
h q[0];
cx q[0], q[1];
\end{verbatim}

\subsection{Registers and timing}

OpenQASM declares \texttt{qubit[n]} quantum and \texttt{bit[m]} classical registers. Statements execute in order unless grouped into \texttt{box} blocks for timing semantics. Gate arguments specify qubit indices; creg stores measurement outcomes for conditional gates such as \texttt{if (c == 1) x q[0];}.

\subsection{Reading and writing QASM}

Export from visual builders or SDKs (Qiskit \texttt{qasm3.dumps}, Cirq \texttt{to\_qasm}). Always map each statement back to unitaries on the Hilbert space; \texttt{cx} is CNOT, \texttt{h} is Hadamard, \texttt{rz(theta)} is $R_z(\theta)$.

\begin{example}
The OpenQASM snippet above corresponds to $U=\mathrm{CNOT}_{0,1}\,(H\otimes I)$ applied left-to-right as $U\ket{00}=\ket{\Phi^+}$. Reversing gate order in the file changes the unitary product accordingly.
\end{example}

\begin{example}
The unitary for $H\otimes I$ on two qubits is
\[
H\otimes I=\frac{1}{\sqrt{2}}\begin{bmatrix}
1&1&0&0\\1&-1&0&0\\0&0&1&1\\0&0&1&-1
\end{bmatrix}.
\]
Apply it to $\ket{00}$ before CNOT to derive~\eqref{eq:bell-circuit}.
\end{example}

\section{Depth and Resource Analysis}
\label{sec:depth}

\subsection{Ripple adders revisited}

An $n$-bit ripple-carry adder built from Toffoli and CNOT has depth $O(n)$ in the worst case---analogous classical depth from Chapter~\ref{ch:classical}. Quantum carry-lookahead and Fourier adders reduce depth at the cost of ancillas.

\subsection{Simulation cost}

Exact state-vector simulation of $n$ qubits requires $O(2^n)$ complex amplitudes. Clifford circuits (generated by $\{H,S,\mathrm{CNOT}\}$ without $T$) simulate in polynomial time (Gottesman--Knill), explaining why random-circuit supremacy demos use non-Clifford gates \cite{arute2019}.

\section{Modern Software Ecosystem (2026)}
\label{sec:software}

Production quantum software typically combines:
\begin{itemize}[noitemsep]
\item \textbf{Circuit construction:} OpenQASM~3, Qiskit, Cirq, PennyLane, and vendor SDKs.
\item \textbf{Simulation:} local state-vector simulators (exponential in $n$) and stabilizer simulators for Clifford circuits.
\item \textbf{Execution:} cloud access to superconducting, trapped-ion, neutral-atom, and photonic backends with calibration data and error models.
\end{itemize}

Map API calls (\texttt{h(0)}, \texttt{cx(0,1)}) to unitaries and tensor structure before trusting black-box outputs.

\section*{Exercises}
\begin{enumerate}[label=\arabic*.]
\item Write the $4\times 4$ matrix for $H\otimes I$ and verify $ (H\otimes I)\ket{00}$.
\item Count depth (qualitatively) of an $n$-bit ripple-carry adder from Toffoli gates.
\item Export a Bell circuit to OpenQASM: which gates appear?
\item Why does CNOT not equal $\mathrm{CNOT}_{1,0}$? Give the matrix of swapped control/target.
\item Estimate memory for simulating $n=30$ qubits in double precision (bytes $\approx 16\cdot 2^n$).
\item Which gates in $\{H,S,\mathrm{CNOT}\}$ generate the Clifford group?
\item Transpile $T$ into $R_z(\pi/4)$ up to global phase.
\item Explain why measurement is not modeled as a unitary gate in~\eqref{eq:circuit-evolution}.
\end{enumerate}
